%=============================================================================
%  Chapitre 9 — Sprint 6 : Portail Agence, Journal d'Audit et Finalisation
%=============================================================================
\chapter{Sprint 6 --- Portail Agence, Journal d'Audit et Finalisation}

\begin{tcolorbox}[sprintbox,title={Sprint 6 --- Portail Agence et Audit}]
\textbf{Durée :} 2 semaines \quad
\textbf{Points de Story :} 37 \quad
\textbf{Objectif :} Livrer le hub conciergerie agence avec streaming en direct des maisons,
un tableau de bord de journal d'audit complet, les demandes de permission résidents,
les modes de sécurité et la gestion du profil.
\end{tcolorbox}

\section{Hub Conciergerie et Portail Agence}

\subsection{Server-Sent Events (SSE) pour le Statut en Direct}

\begin{lstlisting}[style=codeblock,language=JavaScript,caption={Endpoint SSE pour le Hub Conciergerie (backend/routes/concierge.js)}]
router.get('/stream', protect, agencyOnly, (req, res) => {
  res.setHeader('Content-Type',  'text/event-stream');
  res.setHeader('Cache-Control', 'no-cache');
  res.setHeader('Connection',    'keep-alive');
  res.flushHeaders();

  const push = (data) => res.write(`data: ${JSON.stringify(data)}\n\n`);

  // Envoyer l'instantane initial de tous les statuts admin
  push({ type: 'snapshot', houses: getOnlineAdmins() });

  // S'abonner aux evenements de connexion admin
  const onConnect    = (admin) => push({ type:'online',  houseCode: admin.houseCode });
  const onDisconnect = (admin) => push({ type:'offline', houseCode: admin.houseCode });

  adminEvents.on('connect',    onConnect);
  adminEvents.on('disconnect', onDisconnect);

  // Nettoyage a la deconnexion du client
  req.on('close', () => {
    adminEvents.off('connect',    onConnect);
    adminEvents.off('disconnect', onDisconnect);
    res.end();
  });
});
\end{lstlisting}

\section{Système de Journal d'Audit}

\subsection{Schéma AuditLog}

\begin{lstlisting}[style=codeblock,language=JavaScript,caption={Sch\'ema AuditLog (backend/models/AuditLog.js)}]
const auditLogSchema = new mongoose.Schema({
  actor:      { type: mongoose.Schema.Types.ObjectId, ref: 'User' },
  actorName:  { type: String, required: true },
  actorRole:  { type: String },
  action:     { type: String, required: true },
  category:   { type: String,
                enum: ['auth','device','security','user','scenario','system'],
                default: 'system' },
  details:    { type: String },
  metadata:   { type: mongoose.Schema.Types.Mixed },
  ipAddress:  { type: String },
  houseCode:  { type: String }
}, { timestamps: true });
\end{lstlisting}

\subsection{Événements Journalisés}

\begin{center}
\begin{longtable}{p{0.16\textwidth} p{0.72\textwidth}}
\caption{Catégories et Événements du Journal d'Audit}\label{tab:audit_events}\\
\toprule
\textbf{Catégorie} & \textbf{Événements Journalisés} \\
\midrule
\endfirsthead
\toprule
\textbf{Catégorie} & \textbf{Événements Journalisés} \\
\midrule
\endhead
\bottomrule
\endfoot
\texttt{auth}     & Connexion (succès/échec), inscription, réinitialisation mdp, vérification email, 2FA activer/désactiver, connexion Google OAuth \\
\texttt{device}   & Appareil créé, mis à jour, supprimé, commande exécutée (ON/OFF/luminosité/scène) \\
\texttt{security} & Urgence gaz activée, urgence levée, mode Absence activé, mode Verrouillage activé \\
\texttt{user}     & Résident ajouté/supprimé, rôle modifié, permission accordée/révoquée \\
\texttt{scenario} & Scénario créé/mis à jour/supprimé, scénario déclenché par le planificateur \\
\texttt{system}   & Démarrage serveur, démarrage courtier MQTT, initialisation planificateur \\
\end{longtable}
\end{center}

\section{Modes de Sécurité}

\begin{center}
\begin{longtable}{p{0.18\textwidth} p{0.70\textwidth}}
\caption{Définition des Modes de Sécurité}\label{tab:security_modes}\\
\toprule
\textbf{Mode} & \textbf{Comportement} \\
\midrule
\endfirsthead
\toprule
\textbf{Mode} & \textbf{Comportement} \\
\midrule
\endhead
\bottomrule
\endfoot
Normal           & Fonctionnement standard ; tous les appareils librement contrôlables par les utilisateurs autorisés. \\
Mode Absence     & Surveillance caméra renforcée ; toutes les portes extérieures verrouillées ; fenêtres fermées ; vanne gaz fermée ; événements de mouvement journalisés. \\
Mode Verrouillage & Sécurité maximale ; toutes les caméras armées ; toutes les portes et fenêtres verrouillées ; vanne gaz fermée ; contrôle des appareils résidents suspendu. \\
\end{longtable}
\end{center}

\section{Récapitulatif Complet des Routes API}

\begin{center}
\begin{longtable}{p{0.26\textwidth} p{0.14\textwidth} p{0.48\textwidth}}
\caption{Toutes les Routes Backend API}\label{tab:api_summary}\\
\toprule
\textbf{Préfixe de Route} & \textbf{Auth} & \textbf{Rôle} \\
\midrule
\endfirsthead
\toprule
\textbf{Préfixe de Route} & \textbf{Auth} & \textbf{Rôle} \\
\midrule
\endhead
\bottomrule
\endfoot
\api{/api/auth}         & Aucune/JWT   & Inscription, connexion, vérification email, réinitialisation mdp \\
\api{/api/oauth}        & Aucune       & Callback Google OAuth 2.0 \\
\api{/api/2fa}          & JWT          & Configuration, vérification, désactivation TOTP 2FA \\
\api{/api/profile}      & JWT          & Obtenir/mettre à jour le profil et l'avatar utilisateur \\
\api{/api/users}        & Admin JWT    & Liste, mise à jour, suppression des utilisateurs de la maison \\
\api{/api/devices}      & Admin JWT    & CRUD des appareils \\
\api{/api/floorplan}    & JWT          & État plan, commandes, scènes \\
\api{/api/scenarios}    & Admin JWT    & CRUD des scénarios \\
\api{/api/gas}          & Admin JWT    & État gaz, seuil, levée urgence \\
\api{/api/permissions}  & JWT          & CRUD demandes de permission et approbation \\
\api{/api/messages}     & JWT          & Soumission de message support et réponse admin \\
\api{/api/agency}       & Agence JWT   & Liste des admins pour la vue agence \\
\api{/api/units}        & Admin JWT    & Gestion des unités/maisons \\
\api{/api/concierge}    & Agence JWT   & Données hub conciergerie et flux SSE \\
\api{/api/audit}        & Admin JWT    & Requête et pagination du journal d'audit \\
\api{/api/stats}        & Admin JWT    & Statistiques tableau de bord \\
\api{/api/energy}       & Admin JWT    & Consommation énergétique horaire/quotidienne \\
\api{/api/companion}    & JWT          & Endpoint chat assistant IA \\
\api{/api/owner-requests}& Aucune/JWT  & Demandes de compte propriétaire \\
\end{longtable}
\end{center}

\section{Revue et Rétrospective Sprint 6}

\begin{tcolorbox}[goalbox,title={Revue Sprint 6}]
\textbf{Complété :} 8 stories sur 8 (37 PS). Plateforme fonctionnellement complète. \\
\textbf{Points forts :} Le hub conciergerie agence a affiché 3 maisons simulées avec le
clignotement en ligne/hors ligne en direct. Le journal d'audit a été exporté avec succès
en PDF avec jsPDF. Le mode Absence a immédiatement verrouillé tous les servomoteurs de
porte sur l'ESP32. \\
\textbf{Rétrospective :}
\begin{itemize}[leftmargin=1cm,itemsep=2pt]
  \item \textbf{Ce qui a bien fonctionné :} SSE était plus simple que Socket.IO pour le
        cas d'utilisation de streaming en lecture seule de l'agence.
  \item \textbf{Défi :} Le formatage de l'export PDF a nécessité un réglage minutieux des
        largeurs de colonnes dans jsPDF pour les longues chaînes d'action d'audit.
  \item \textbf{Amélioration :} Ajout d'une limite de 2~Mo sur les uploads d'avatar après
        qu'une image base64 de 15~Mo a causé des avertissements de taille de document MongoDB.
\end{itemize}
\end{tcolorbox}
